\chapter{Process Identification}

\section{Theoretical Background}

\subsection{First-Order Plus Time Delay Model}

A large class of industrial processes --- including thermal systems --- can be adequately
described, in the neighbourhood of an operating point, by a First-Order Plus Time Delay
(FOPTD) transfer function:

\begin{equation}
  \boxed{G(s) = \frac{K \cdot e^{-\theta s}}{\tau s + 1}},
  \label{eq:foptd_id}
\end{equation}

where $K$ is the static gain, $\tau$ [\si{\second}] the time constant, and $\theta$
[\si{\second}] the dead time. For the inner process of this installation:

\begin{itemize}[leftmargin=*, itemsep=3pt]
  \item \textbf{Input:} $F_{1,SP}$ [\si{\metre\cubed\per\hour}]
  \item \textbf{Output:} $T_{in1,HE1}$ [\si{\degreeCelsius}]
  \item \textbf{$K$:} $[\si{\degreeCelsius}/(\si{\metre\cubed\per\hour})]$ ---
        static temperature change per unit flow step
\end{itemize}

\textbf{Justification.} The furnace acts as a thermally well-mixed chamber: energy input
is proportional to power, and heat removal is proportional to the outgoing flow. This
single-dominant-lag structure is well approximated by a first-order system. The dead time
$\theta$ captures the transport delay as fluid travels from the valve to the temperature
sensor. Higher-order effects (sensor dynamics, secondary thermal masses) are absorbed into
the effective $\tau$ and $\theta$ values.

\subsection{Step Response Characterisation}

Given a step input $\Delta F_{1,SP}$ applied at $t = t_0$, the ideal FOPTD step response is:
\begin{equation}
  \Delta T_{in1}(t) =
  \begin{cases}
    0 & t < t_0 + \theta \\
    K \cdot \Delta F_{1,SP}
      \left(1 - e^{-(t - t_0 - \theta)/\tau}\right) & t \geq t_0 + \theta
  \end{cases}\,.
  \label{eq:foptd_step}
\end{equation}

The three parameters are extracted graphically from the experimental curve:

\begin{align}
  K   &= \frac{\Delta T_{in1,\infty}}{\Delta F_{1,SP}},
  \label{eq:K_id}\\[4pt]
  \tau &: \quad T_{in1}(t_{63}) = T_{in1,\text{base}} + 0.632\;\Delta T_{in1,\infty}
  \quad \Rightarrow \quad \tau = t_{63} - (t_0 + \theta),
  \label{eq:tau_id}\\[4pt]
  \theta &\approx t_{\text{first}} - t_0
  \quad \text{where } |\Delta T_{in1}| > 0.02\;\Delta T_{in1,\infty}.
  \label{eq:theta_id}
\end{align}

Figure~\ref{fig:foptd_response} illustrates the ideal step response with these
characteristic points annotated.

\begin{figure}[H]
  \centering
  \begin{tikzpicture}
    \begin{axis}[
      width=12cm, height=6cm,
      xlabel={Time [\si{\second}]},
      ylabel={$\Delta T_{in1}$ [\si{\degreeCelsius}]},
      xmin=0, xmax=240,
      ymin=-0.5, ymax=12,
      xtick={30,45,90,150,210},
      xticklabels={$t_0$, $t_0{+}\theta$, $t_{63}$, , },
      ytick={0,7.57,12},
      yticklabels={$0$, $0.632\,\Delta T_\infty$, $\Delta T_\infty$},
      grid=both,
      grid style={line width=0.3pt, draw=gray!30},
      major grid style={line width=0.5pt, draw=gray!50},
      axis lines=left,
      tick label style={font=\small},
      label style={font=\small},
    ]
      % FOPTD step response: K=12, tau=45s, theta=15s, step at t=30s
      \addplot[domain=0:30, thick, draw=darkblue] {0};
      \addplot[domain=30:45, thick, draw=darkblue] {0};
      \addplot[domain=45:240, thick, draw=darkblue, samples=80]
        {12*(1 - exp(-(x-45)/45))};

      % Asymptote
      \addplot[dashed, thin, draw=gray] coordinates {(0,12)(240,12)};

      % 63.2% line
      \addplot[dashed, thin, draw=orange!80!black]
        coordinates {(0,7.57)(90,7.57)(90,0)};

      % Dead time arrow
      \draw[-Stealth, thick, draw=red!70!black]
        (axis cs:30,1.5) -- (axis cs:45,1.5)
        node[midway, above, font=\scriptsize, red!70!black]{$\theta$};

      % Time constant arrow
      \draw[<->, thick, draw=blue!70!black]
        (axis cs:45,-0.3) -- (axis cs:90,-0.3)
        node[midway, below, font=\scriptsize, blue!70!black]{$\tau$};

      \node[font=\scriptsize, right] at (axis cs:210,11.6) {$K\,\Delta F_{1,SP}$};
    \end{axis}
  \end{tikzpicture}
  \caption{Ideal FOPTD step response with characteristic parameters $K$, $\tau$, $\theta$
           annotated (illustrative values: $K=12$\,\si{\degreeCelsius\per\metre\cubed\per\hour},
           $\tau=\SI{45}{\second}$, $\theta=\SI{15}{\second}$).}
  \label{fig:foptd_response}
\end{figure}

\subsection{IMC Tuning Rules}
\label{sec:imc}

Internal Model Control (IMC) provides an explicit relationship between the desired
closed-loop speed of response and the PI controller gains. For a FOPTD process,
the IMC-derived PI controller has the following gains~\cite{Rivera1986}:

\begin{equation}
  \boxed{K_p = \frac{\tau}{K\,(\lambda + \theta)}, \qquad T_i = \tau},
  \label{eq:imc}
\end{equation}

where $\lambda$ [\si{\second}] is the desired closed-loop time constant --- the single
free tuning parameter. Three pre-defined levels are provided:

\begin{table}[H]
  \centering
  \caption{IMC tuning levels and corresponding $\lambda$ values.}
  \label{tab:imc_levels}
  \begin{tabular}{@{}L{3.5cm}C{3.0cm}L{6.5cm}@{}}
    \toprule
    \textbf{Level} & \textbf{$\lambda$} & \textbf{Expected behaviour} \\
    \midrule
    Aggressive  & $0.5\,\tau$   & Fast response, $\sim$10\,\% overshoot possible \\
    Normal      & $\tau$        & Balanced --- recommended starting point \\
    Conservative& $2\,\tau$     & Slow, robust --- use if oscillations appear \\
    \bottomrule
  \end{tabular}
\end{table}

\textit{Note:} $\lambda \geq \theta$ is a physical lower bound (closed-loop time constant
cannot be shorter than the process dead time for a PI controller). The implementation
enforces $\lambda \leftarrow \max(\lambda, \theta)$.

\textbf{Advantage over Ziegler--Nichols:} IMC rules have a clear physical interpretation
--- $\lambda$ directly sets the desired closed-loop response speed. Detuning is trivial:
simply increase $\lambda$.

\section{Experimental Protocol}

\subsection{Pre-Test Conditions}

Before starting the step test, the following conditions must be verified:

\begin{enumerate}[leftmargin=*, itemsep=3pt]
  \item The furnace has been running for at least 30\,minutes; process is at steady state
    (stability criterion: $|T_{in1}|$ drift $<\SI{0.5}{\degreeCelsius}$ over
    5\,minutes).
  \item The cascade controller is in \textbf{MANUAL} mode. The identification API call
    is rejected if mode~$\neq$~MANUAL.
  \item Baseline values recorded: $T_{hin}$, $T_{in1,HE1}$, $T_{out2,HE1}$, $F_1$, $F_2$.
\end{enumerate}

\subsection{Step Amplitude and Duration}

The step amplitude should be 10--20\,\% of the $F_1$ operating range.

\begin{itemize}[leftmargin=*, itemsep=3pt]
  \item \textbf{Too small:} temperature rise buried in sensor noise
    (typical resolution $\sim\SI{0.5}{\degreeCelsius}$).
  \item \textbf{Too large:} nonlinear region; identified parameters not representative of
    the linearised operating point.
  \item \textbf{Recommended:} $\Delta F_{1,SP} = +\SI{0.5}{\metre\cubed\per\hour}$
    from nominal $\approx \SI{3.0}{\metre\cubed\per\hour}$ ($\approx 10\,\%$ of range).
\end{itemize}

A duration of \SI{180}{\second} is chosen as a default ($\approx 3$--$4\,\tau$ for an
estimated $\tau$ of 30--60\,s). The test can be cancelled early if the transient is
clearly complete.

\subsection{Acquisition and Identification Software}

The \code{StepIdentifier} class in \code{identification.py} orchestrates the test
through three phases:

\begin{description}[leftmargin=2.5cm, style=nextline, itemsep=4pt]
  \item[PRE\_STEP] (5 samples $\times$ \SI{3}{\second} = 15\,\si{\second}):
    The identifier records the baseline. $F_{1,SP}$ is held at its current value.
    Baseline $T_{in1}$ is the average of the 5 pre-step samples.

  \item[STEP] (configurable duration):
    At the first poll cycle after PRE\_STEP completes, $F_{1,SP}$ is set to
    $F_{1,\text{base}} + \text{amplitude}$. All (timestamp, $T_{in1}$, $T_{out2}$,
    $F_{1,SP}$) samples are recorded. Progress is broadcast via \code{ident\_update}.

  \item[DONE] (automatic):
    \code{\_analyze()} is called, parameters are extracted, $F_{1,SP}$ is restored,
    IMC gains are computed and returned to the web UI.
\end{description}

\textbf{Sampling resolution.} With $T_s = \SI{3}{\second}$, the resolution of $\tau$ and
$\theta$ is $\pm\SI{3}{\second}$, with linear interpolation between samples improving it
to $\pm\SI{1.5}{\second}$ in practice. For thermal processes where $\tau \gg \theta$, this
is acceptable for PI tuning purposes.

\textbf{Safety.} If the test is cancelled via \code{POST /api/identification/cancel},
$F_{1,SP}$ is immediately restored to the baseline value. The cascade mode cannot be
changed during a running test.

\section{Identification Results}

\pendingbox{%
  Results from the experimental session of 2026-05-19 with Michal. Complete
  Section~5.3.1 to 5.3.4 after the session.%
}

\subsection{Raw Step Response Curves}

\begin{figure}[H]
  \centering
  \fbox{\begin{minipage}{0.85\textwidth}
    \centering\vspace{2cm}
    \textit{[Figure: Step response curves exported from the \texttt{/cascade}
    identification chart.\\
    Left axis: $T_{in1,HE1}$ (red), $T_{out2,HE1}$ (orange dashed).\\
    Right axis: $F_{1,SP}$ step (blue).\\
    X-axis: time in \si{\second} from step application.]}
    \vspace{2cm}
  \end{minipage}}
  \caption{Measured step response of $T_{in1,HE1}$ to a $\Delta F_{1,SP}$
    step~--- to be completed.}
  \label{fig:step_response}
\end{figure}

\subsection{Extracted Parameters}

\begin{table}[H]
  \centering
  \caption{Identified FOPTD parameters.}
  \label{tab:id_results}
  \begin{tabular}{@{}L{5cm}C{3cm}C{2cm}L{3cm}@{}}
    \toprule
    \textbf{Parameter} & \textbf{Value} & \textbf{Unit} & \textbf{Method} \\
    \midrule
    $\Delta F_{1,SP}$ (step amplitude)   & TBD & \si{\metre\cubed\per\hour} & Imposed \\
    $\Delta T_{in1,\infty}$ (steady state)& TBD & \si{\degreeCelsius} & Measured \\
    $K$ (static gain)                    & TBD & \si{\degreeCelsius\per\metre\cubed\per\hour} & Eq.~\eqref{eq:K_id} \\
    $\tau$ (time constant)               & TBD & \si{\second} & 63.2\% crossing \\
    $\theta$ (dead time)                 & TBD & \si{\second} & 2\% threshold \\
    \bottomrule
  \end{tabular}
\end{table}

\subsection{IMC Gain Calculations}

\begin{table}[H]
  \centering
  \caption{IMC-tuned PI gains at three levels (values pending experimental results).}
  \label{tab:imc_results}
  \begin{tabular}{@{}L{3.5cm}C{2.5cm}C{2.5cm}C{2.5cm}@{}}
    \toprule
    \textbf{Tuning level} & \textbf{$\lambda$ [\si{\second}]} &
    \textbf{$K_p$} & \textbf{$T_i$ [\si{\second}]} \\
    \midrule
    Aggressive  & $0.5\,\tau$ & TBD & TBD \\
    Normal      & $\tau$      & TBD & TBD \\
    Conservative& $2\,\tau$   & TBD & TBD \\
    \bottomrule
  \end{tabular}
\end{table}

\subsection{Model Validation}

\pendingbox{Comparison of identified FOPTD model prediction vs.\ measured step response.
  Report RMSE and discuss validity domain.}
